\documentclass[10pt,a4paper]{amsart}
\usepackage[margin=19mm]{geometry}
\usepackage{amsmath,amssymb,mathtools}
\usepackage[scr=rsfs,cal=euler]{mathalfa}
\usepackage{xfrac}
\usepackage[unicode,hidelinks,bookmarks=false]{hyperref}
\newcommand{\ie}{i.e.\ }
\newcommand{\cf}{cf.\ }
\newcommand{\resp}{respectively\ }
\newcommand{\Subsection}{\subsection}
\providecommand{\nobreakdash}{\nobreak\hbox{-}}
\setlength{\emergencystretch}{2em}
\multlinegap0pt
\usepackage{amssymb}
\usepackage{csquotes}
\newtheorem{theorem}{Theorem}
\title{Integers representable as a difference of two rational fourth powers}
\author{Ashleigh Ratcliffe, Nguyen Xuan Tho}
\date{Source: arXiv:2604.15832v2 (CC BY 4.0)}
\begin{document}
\maketitle

\noindent\emph{Source and licence.} Integers representable as a difference of two rational fourth powers. Version arXiv:2604.15832v2 (CC BY 4.0), distributed by arXiv under the Creative Commons Attribution 4.0 International licence, http://creativecommons.org/licenses/by/4.0/, as recorded in the arXiv licence metadata for this exact version and shown on the abstract page https://arxiv.org/abs/2604.15832v2. Full licence notice: \texttt{licenses/arxiv-2604.15832-CC-BY-4.0.txt}.\par\smallskip

\noindent\textit{Editorial scope.} Counted unit: the article's theorem T2 (source0.tex lines 116--125). The article's complete original proof of it is delivered with it (source0.tex lines 128--157, one \texttt{proof} environment of 30 lines). No mathematics is written, repaired or re-derived: the statement and the proof are the article's own lines, in the article's own notation, and no retained line is substituted or re-flowed.  Retained uncounted as the source-local context the proof needs: lines 80--82.  Nothing was omitted from any retained span, so the package carries no omission record.  The retained lines cite 1 works, whose entries are transcribed verbatim from the article's own bibliography source in the e-print and delivered with short editorial labels recorded in the item's reference mapping.  No retained line contains a figure environment, an image-inclusion command or drawing code, so no image is delivered and the package declares no asset dependency.  Editorial formatting only, in addition: an AMS \texttt{amsart} preamble that restates the article's own declarations and macros used by the retained lines, namely the environments \texttt{theorem} on separate counters, together with the standard packages those lines need.  The retained environments print in this document as Theorem 1 in order of appearance, whereas the article prints the same items with the numbering of its own full document.  The complete native source of the article as distributed in the arXiv e-print for this exact version is bundled unchanged as \texttt{sources/198643/source0.tex}.
	\begin{equation}\label{eq:x4my4mnz4}
		x^4-y^4=nz^4
	\end{equation} 

\begin{theorem}\label{T2}\cite[Section 6.3.4]{mainbook} 
We can reduce the problem of solving \eqref{eq:x4my4mnz4} in integers to solving equations of the form
\begin{equation}\label{eq:abcuvw}
a^2u^8+b^2v^8=c w^4
\end{equation}
for pairwise coprime positive integers $a,b,c$ with $abc\in\{n/8,2n\}$ and integer variables $u,v,w$ such that
\begin{equation}\label{eq:cond}
\gcd(au,bv)= \gcd(au,cw)= \gcd(bv,cw)= 1.
\end{equation}
\end{theorem}

\begin{proof} 
To be self-contained, we include the proof. Assume that an integer solution $(x,y,z)$ with $xyz \neq 0$ to \eqref{eq:x4my4mnz4} exists. Without loss of generality, we can assume that $x,y,z >0$, since if $(x,y,z)$ is a solution, then so are $(\pm x,\pm y, \pm z)$. We may further assume that $\gcd(x,y,z)=1$, as if $d=\gcd(x,y,z)>1$, then $(x/d,y/d,z/d)$ is also a solution with $\gcd(x/d,y/d,z/d)=1$.  We can further assume that $x,y,z$ are pairwise coprime. If $x$ and $y$ share a common factor $p$, then the left-hand side of \eqref{eq:x4my4mnz4} is divisible by $p^4$, and therefore the right-hand side must also be divisible by $p^4$, and $n$ is not divisible by a fourth power, so $z$ must be. A similar argument shows that $\gcd(x,z)=\gcd(y,z)=1$. 
Therefore, either $x$ and $y$ are both odd, or $x$ and $y$ have different parities. 

We first consider the case that $x$ and $y$ are both odd. We use the change of variables $t=(x+y)/2$, $s=(x-y)/2$, where $t$ and $s$ have different parity and are coprime. After substituting into \eqref{eq:x4my4mnz4}, we obtain
$$
8st(s^2+t^2)=nz^4.
$$
We must have either that $n$ is divisible by $8$, in which case let $n'=n/8$ and $z=r$, otherwise, $z$ is even so can be written as $z=2r$ for integer $r$ and let $n'=2n$. In either case, we obtain 
\begin{equation}\label{eq:red}
st(s^2+t^2)=n'r^4.
\end{equation}

Next we consider the case that $x$ and $y$ have different parity. We use the change of variables $t=x+y$ and $s=x-y$, where $s$ and $t$ are both odd and coprime (because $x$ and $y$ are coprime). After substituting into \eqref{eq:x4my4mnz4}, we obtain
$$
st(s^2+t^2)=2nz^4,
$$
letting $n'=2n$ and $z=r$, we obtain \eqref{eq:red}.
%We must have either that $n$ is divisible by $8$, in which case let $n'=n/8$ and $z=r$, otherwise, $z$ is even so can be written as $z=2r$ for integer $r$, and we obtain \eqref{eq:red} where $n'=2n$. 

Because $s$ and $t$ are coprime, $s^2+t^2$ is also coprime to $s$ and $t$. Solutions to \eqref{eq:red} must satisfy
\begin{equation}\label{stabc}
s=au^4, \quad t=bv^4, \quad s^2+t^2=c w^4,
\end{equation}
where $\gcd(au,bv)=\gcd(au,cw)=\gcd(bv,cw)=1$ and $abc=n'$.
After substituting \eqref{stabc} into \eqref{eq:red}, we obtain 
$$
a^2u^8+b^2v^8=c w^4.
$$
\end{proof}

\begin{thebibliography}{99}
\bibitem[mainbo]{mainbook} Bogdan Grechuk. ``Polynomial Diophantine Equations. A Systematic Approach.'' \emph{Springer} (2024). %
\end{thebibliography}
\end{document}
